% PART 3: UC16 - UC22
\section*{1.11. Đặc tả Use Case Scenario (Phần 4: UC 16--20)}

\begin{longtable}{|L{3.5cm}|L{12.5cm}|}
\hline
\textbf{Use Case ID} & UC16 \\ \hline
\textbf{Use Case Name} & Đánh dấu món ăn đã hoàn tất và ``Sẵn sàng phục vụ'' (Ready to Serve) \\ \hline
\textbf{Primary Actors} & Đầu bếp (Kitchen Staff) \\ \hline
\textbf{Description} & Ghi nhận món ăn đã nấu xong, báo cho hệ thống để điều phối nhân viên phục vụ mang món ra bàn. \\ \hline
\textbf{Trigger} & Đầu bếp hoàn tất việc chế biến và bày món ăn ra khay chờ (Expo window). \\ \hline
\textbf{Preconditions} & $\bullet$ Món ăn đang ở trạng thái ``Đang chế biến'' (Cooking) trên màn hình KDS. \\ \hline
\textbf{Postconditions} &
$\bullet$ Trạng thái món ăn đổi thành ``Sẵn sàng phục vụ'' (Ready to Serve).\newline
$\bullet$ Thông báo được chuyển đến màn hình POS/thiết bị của nhân viên phục vụ. \\ \hline
\textbf{Normal Flow} &
1. Đầu bếp hoàn tất việc chế biến và bày món ăn ra khay chờ (Expo window).\newline
2. Đầu bếp chạm/bấm nút ``Hoàn thành'' (Done/Bump) trên phiếu món ăn tương ứng tại KDS.\newline
3. Hệ thống xóa phiếu món ăn đó khỏi màn hình của Đầu bếp (hoặc chuyển sang danh sách Đã xong).\newline
4. Hệ thống cập nhật trạng thái món thành ``Sẵn sàng phục vụ'' trên toàn hệ thống.\newline
5. Hệ thống gửi thông báo nhấp nháy hoặc âm thanh đến POS của nhân viên phục vụ phụ trách bàn đó. \\ \hline
\textbf{Alternative Flows} &
\textbf{2a. Bấm nhầm nút khi chưa nấu xong:} Đầu bếp truy cập vào tab ``Đã hoàn thành'' (History/Done), chọn phiếu vừa bấm và nhấn ``Hoàn tác'' (Recall). Hệ thống đưa phiếu trở lại màn hình chính đang nấu. \\ \hline
\end{longtable}

\begin{longtable}{|L{3.5cm}|L{12.5cm}|}
\hline
\textbf{Use Case ID} & UC17 \\ \hline
\textbf{Use Case Name} & Xem cảnh báo các món ăn sắp trễ thời gian phục vụ quy định \\ \hline
\textbf{Primary Actors} & Đầu bếp (Kitchen Staff) \\ \hline
\textbf{Description} & Hệ thống hỗ trợ đầu bếp phân bổ lại sự tập trung vào các món sắp hoặc đã vượt quá thời gian chế biến tiêu chuẩn (SLA). \\ \hline
\textbf{Trigger} & Hệ thống phát hiện một phiếu nấu đạt đến 80\% thời gian chế biến cho phép. \\ \hline
\textbf{Preconditions} &
$\bullet$ Hệ thống đang đếm ngược thời gian nấu của ít nhất một món ăn.\newline
$\bullet$ KDS đang hoạt động bình thường. \\ \hline
\textbf{Postconditions} & $\bullet$ Đầu bếp nhận biết được sự khẩn cấp và ưu tiên xử lý món bị trễ. \\ \hline
\textbf{Normal Flow} &
1. Hệ thống liên tục so sánh thời gian đã trôi qua của một phiếu nấu với Thời gian chuẩn bị dự kiến (Prep Time) trong cơ sở dữ liệu.\newline
2. Khi phiếu đạt đến 80\% thời gian cho phép, hệ thống chuyển viền phiếu sang màu vàng (Cảnh báo sắp trễ).\newline
3. Khi phiếu vượt quá 100\% thời gian, hệ thống chuyển toàn bộ nền phiếu sang màu đỏ và nhấp nháy liên tục.\newline
4. Đầu bếp nhìn thấy cảnh báo trực quan trên KDS và điều chỉnh ưu tiên nấu nướng.\newline
5. (Tuỳ chọn) Hệ thống tự động gửi cảnh báo cho Quản lý bếp qua màn hình tổng. \\ \hline
\textbf{Alternative Flows} &
\textbf{2a. Món ăn bị đình trệ do thiếu nguyên liệu:} Đầu bếp nhấn vào phiếu và chọn ``Tạm ngưng'' (Hold) kèm lý do (Hết nguyên liệu). Hệ thống dừng đếm thời gian trễ và báo lại cho quản lý/nhân viên phục vụ. \\ \hline
\end{longtable}

\begin{longtable}{|L{3.5cm}|L{12.5cm}|}
\hline
\textbf{Use Case ID} & UC18 \\ \hline
\textbf{Use Case Name} & Cập nhật thủ công số lượng nguyên liệu đã sử dụng sau khi nấu xong \\ \hline
\textbf{Primary Actors} & Đầu bếp (Kitchen Staff) \\ \hline
\textbf{Description} & Trừ đi lượng nguyên liệu bổ sung, hư hỏng hoặc làm rơi vãi không theo định lượng chuẩn, giúp số liệu tồn kho luôn chính xác. \\ \hline
\textbf{Trigger} & Đầu bếp phát hiện hao hụt nguyên liệu ngoài định mức và mở chức năng ``Cập nhật tồn kho''. \\ \hline
\textbf{Preconditions} & $\bullet$ Đầu bếp có quyền truy cập vào module quản lý nguyên liệu/tồn kho tại bếp. \\ \hline
\textbf{Postconditions} &
$\bullet$ Cơ sở dữ liệu tồn kho được trừ số lượng tương ứng.\newline
$\bullet$ Lịch sử trừ kho được ghi nhận cho tài khoản Đầu bếp. \\ \hline
\textbf{Normal Flow} &
1. Đầu bếp mở chức năng ``Cập nhật tồn kho'' trên màn hình quản lý.\newline
2. Đầu bếp tìm kiếm tên nguyên liệu (ví dụ: Cá hồi fillet, Dầu ăn).\newline
3. Hệ thống hiển thị số lượng tồn khả dụng hiện tại.\newline
4. Đầu bếp nhập số lượng cần trừ đi (báo hỏng, hao hụt) và ghi chú lý do (VD: ``Làm rớt'').\newline
5. Hệ thống xác nhận giao dịch, cập nhật lại số tồn thực tế và lưu vào nhật ký (audit log). \\ \hline
\textbf{Alternative Flows} &
\textbf{4a. Số lượng trừ vượt quá tồn kho hiện tại:} Hệ thống cảnh báo ``Số lượng trừ không hợp lệ (Vượt quá tồn kho)''. Đầu bếp phải nhập lại số chính xác hoặc yêu cầu quản lý kiểm kê kho vật lý. \\ \hline
\end{longtable}

\begin{longtable}{|L{3.5cm}|L{12.5cm}|}
\hline
\textbf{Use Case ID} & UC19 \\ \hline
\textbf{Use Case Name} & Xuất hóa đơn tổng hợp phí thức ăn, đồ uống và các phụ phí \\ \hline
\textbf{Primary Actors} & Thu ngân (Cashiers) \\ \hline
\textbf{Description} & Hệ thống tự động tính toán tổng số tiền khách hàng cần thanh toán, bao gồm mọi loại phí, để tiến hành thu tiền. \\ \hline
\textbf{Trigger} & Thu ngân chọn bàn yêu cầu tính tiền trên POS và nhấn ``Thanh toán'' (Checkout). \\ \hline
\textbf{Preconditions} &
$\bullet$ Bàn đã ăn xong, khách yêu cầu tính tiền.\newline
$\bullet$ Toàn bộ món ăn trong đơn đã được đánh dấu ``Đã phục vụ'' (Served). \\ \hline
\textbf{Postconditions} & $\bullet$ Hóa đơn tạm tính (Pre-receipt) được in ra hoặc hiển thị thành công với con số chính xác tuyệt đối. \\ \hline
\textbf{Normal Flow} &
1. Thu ngân chọn bàn yêu cầu tính tiền trên POS và nhấn ``Thanh toán'' (Checkout).\newline
2. Hệ thống tổng hợp tất cả các món ăn đã gọi và giá tiền tương ứng từ cơ sở dữ liệu.\newline
3. Hệ thống tự động tính toán các phụ phí: Thuế VAT (VD: 8\% hoặc 10\%), Phí phục vụ (Service charge).\newline
4. Hệ thống hiển thị màn hình hóa đơn chi tiết với Tạm tính, Tổng phí và Thành tiền cuối cùng.\newline
5. Thu ngân nhấn ``In hóa đơn tạm tính''. Hệ thống điều khiển máy in nhiệt xuất ra giấy để mang ra cho khách. \\ \hline
\textbf{Alternative Flows} &
\textbf{2a. Có món ăn đang kẹt ở trạng thái ``Đang chế biến'':} Hệ thống cảnh báo ``Bàn này vẫn còn món chưa phục vụ. Bạn có chắc muốn tính tiền?''. Thu ngân phải kiểm tra lại với phục vụ để hủy món hoặc chờ phục vụ nốt. \\ \hline
\end{longtable}

\begin{longtable}{|L{3.5cm}|L{12.5cm}|}
\hline
\textbf{Use Case ID} & UC20 \\ \hline
\textbf{Use Case Name} & Áp dụng mã khuyến mãi hoặc chương trình giảm giá vào hóa đơn \\ \hline
\textbf{Primary Actors} & Thu ngân (Cashiers) \\ \hline
\textbf{Description} & Hệ thống xác thực và giảm trừ số tiền thanh toán của khách dựa trên mã voucher, hạng thành viên hoặc chương trình khuyến mãi hiện hành. \\ \hline
\textbf{Trigger} & Thu ngân chọn mục ``Khuyến mãi / Giảm giá'' tại màn hình Thanh toán. \\ \hline
\textbf{Preconditions} &
$\bullet$ Đơn hàng đang ở màn hình Thanh toán.\newline
$\bullet$ Chưa chốt phương thức thanh toán cuối cùng (chưa đóng bill). \\ \hline
\textbf{Postconditions} & $\bullet$ Hóa đơn được cập nhật lại số tiền phải trả sau khi áp dụng giảm giá hợp lệ. \\ \hline
\textbf{Normal Flow} &
1. Tại màn hình Thanh toán, Thu ngân chọn mục ``Khuyến mãi / Giảm giá''.\newline
2. Thu ngân quét mã QR voucher, nhập mã code bằng tay hoặc quét thẻ thành viên của khách.\newline
3. Hệ thống kết nối với dịch vụ Khuyến mãi (Promotion Module) để xác thực tính hợp lệ của mã.\newline
4. Hệ thống tính toán số tiền được giảm (VD: giảm 10\% trên tổng bill, tối đa 50k).\newline
5. Hệ thống hiển thị dòng ``Chiết khấu khuyến mãi'' có giá trị âm trên hóa đơn và cập nhật Thành tiền mới.\newline
6. Thu ngân thông báo cho khách số tiền cuối cùng. \\ \hline
\textbf{Alternative Flows} &
\textbf{3a. Mã không hợp lệ hoặc hết hạn:} Hệ thống báo lỗi đỏ: ``Mã khuyến mãi không tồn tại, đã hết hạn hoặc không đủ điều kiện''. Thu ngân thông báo lại cho khách và xóa mã khỏi hóa đơn.\newline
\textbf{3b. Mã giảm giá trùng lặp:} Hệ thống báo lỗi và yêu cầu Thu ngân chỉ chọn 1 mã có lợi nhất cho khách. \\ \hline
\end{longtable}

\section*{1.12. Đặc tả Use Case Scenario (Phần 5: UC 21--25)}

\begin{longtable}{|L{3.5cm}|L{12.5cm}|}
\hline
\textbf{Use Case ID} & UC21 \\ \hline
\textbf{Use Case Name} & Thực hiện chia hóa đơn (Split-bill) theo yêu cầu của khách hàng \\ \hline
\textbf{Primary Actors} & Thu ngân (Cashiers) \\ \hline
\textbf{Description} & Hệ thống tự động phân chia tổng tiền thành nhiều hóa đơn nhỏ để khách hàng dễ dàng trả tiền theo phần ăn của mỗi người hoặc cưa đều. \\ \hline
\textbf{Trigger} & Bàn đã ăn xong và khách yêu cầu chia bill; Thu ngân chọn ``Chia hóa đơn'' (Split Bill). \\ \hline
\textbf{Preconditions} &
$\bullet$ Bàn đã ăn xong, khách yêu cầu chia bill.\newline
$\bullet$ Hóa đơn tổng chưa được thanh toán thành công. \\ \hline
\textbf{Postconditions} & $\bullet$ Hệ thống tách thành nhiều hóa đơn con, mỗi hóa đơn có thể thanh toán bằng các phương thức khác nhau. \\ \hline
\textbf{Normal Flow} &
1. Tại màn hình Thanh toán, Thu ngân chọn ``Chia hóa đơn'' (Split Bill).\newline
2. Hệ thống cung cấp hai lựa chọn: ``Chia đều'' (Split evenly) hoặc ``Chia theo món'' (Split by items).\newline
3. Thu ngân chọn phương thức (ví dụ: Chia theo món).\newline
4. Hệ thống mở giao diện kéo-thả (drag-and-drop). Thu ngân kéo các món ăn sang ``Hóa đơn khách A'', ``Hóa đơn khách B''.\newline
5. Hệ thống tự động tính lại Thuế và Phí phục vụ cho từng hóa đơn con.\newline
6. Thu ngân tiến hành thanh toán tuần tự cho từng hóa đơn nhỏ. \\ \hline
\textbf{Alternative Flows} &
\textbf{5a. Món ăn số lượng 1 nhưng 2 khách muốn cưa đôi:} Thu ngân chọn món đó và bấm ``Chia tỉ lệ'' (Split by fraction). Hệ thống cắt món thành 2 phần (mỗi phần 50\% giá trị) để đưa vào 2 hóa đơn con. \\ \hline
\end{longtable}

\begin{longtable}{|L{3.5cm}|L{12.5cm}|}
\hline
\textbf{Use Case ID} & UC22 \\ \hline
\textbf{Use Case Name} & Xử lý thanh toán qua tiền mặt, thẻ ngân hàng hoặc ví điện tử \\ \hline
\textbf{Primary Actors} & Thu ngân (Cashiers) \\ \hline
\textbf{Description} & Ghi nhận phương thức khách hàng dùng để trả tiền, kết nối với bên thứ 3 (nếu cần) và hoàn tất vòng đời của đơn hàng. \\ \hline
\textbf{Trigger} & Số tiền cuối cùng trên hóa đơn đã được chốt và Thu ngân hỏi khách phương thức thanh toán. \\ \hline
\textbf{Preconditions} & $\bullet$ Số tiền cuối cùng trên hóa đơn đã được chốt. \\ \hline
\textbf{Postconditions} &
$\bullet$ Giao dịch được ghi nhận ``Thành công''.\newline
$\bullet$ Đơn hàng chuyển sang trạng thái ``Đã thanh toán'' (Paid).\newline
$\bullet$ Bàn được giải phóng. \\ \hline
\textbf{Normal Flow} &
1. Thu ngân hỏi khách phương thức thanh toán và chọn tương ứng trên POS (Tiền mặt / Thẻ / Ví điện tử).\newline
2. \textbf{Nếu Tiền mặt:} Thu ngân nhập số tiền khách đưa. Hệ thống tự động tính số tiền thối lại (Change due). Thu ngân bấm ``Hoàn tất''.\newline
3. \textbf{Nếu Thẻ/Ví điện tử:} Hệ thống truyền số tiền sang máy POS quẹt thẻ hoặc màn hình hiển thị mã QR động.\newline
4. Khách hàng quẹt thẻ hoặc quét mã QR.\newline
5. Hệ thống nhận tín hiệu ``Thành công'' từ Payment Gateway trả về.\newline
6. Hệ thống đổi trạng thái đơn thành ``Đã thanh toán'' và in hóa đơn chính thức (VAT/Receipt). \\ \hline
\textbf{Alternative Flows} &
\textbf{5a. Thẻ bị từ chối hoặc quét QR lỗi:} Hệ thống hiển thị popup đỏ ``Giao dịch thất bại'' kèm mã lỗi từ ngân hàng (VD: Số dư không đủ). Thu ngân thông báo khách và hệ thống giữ nguyên hóa đơn để chờ đổi phương thức. \\ \hline
\end{longtable}
